\chapter{Relationship to other Technics instruments}

The message engine described in this reference is not unique to the SX-WSA1R. At least two
other Technics instruments implement the same protocol: the SX-KN5000 and the SX-KN1500.

\section{What is shared}
All three use the \bytes{50} identifier, the same message families, the same short control
messages, the same payload encoding and the same checksum rule. The fixed messages are
byte-for-byte identical:

\begin{center}
\bytes{F0 50 23 7E F7}\quad
\bytes{F0 50 24 7E F7}\quad
\bytes{F0 50 27 7E F7}\quad
\bytes{F0 50 28 7E F7}\quad
\bytes{F0 50 29 7E F7}\quad
\bytes{F0 50 2A 7E F7}
\end{center}

The error mapping is also nearly common: of the thirty-four status values, thirty-two
produce the same report on the SX-WSA1R and the SX-KN5000.

\section{What distinguishes the products}
\textsc{pc}, \textsc{md} and \textsc{ver} together are the discriminator. A message whose
three bytes do not match the receiving instrument is refused with \textsc{error 40}.

\begin{center}
\begin{tabular}{ll}
\toprule
instrument & \textsc{pc} \textsc{md} \textsc{ver} \\
\midrule
SX-WSA1 / SX-WSA1R & \bytes{04 00 11}, \bytes{04 01 11} \\
SX-KN5000 & \bytes{01 28 12} \\
SX-KN1500 & \bytes{01 24 11}, \bytes{01 25 11}, \bytes{01 26 11} \\
\bottomrule
\end{tabular}
\end{center}

The split is legible once the fields are named: \textsc{pc}~=~\bytes{04} is the
\textsc{wsa} product category and \bytes{01} the keyboards, \textsc{md} separates models
within a category, and \textsc{ver} is the exclusive specification's own version --- 2.1
for the WSA1 and the KN1500, 2.2 for the KN5000.

A librarian that supports more than one of these instruments needs only to vary those bytes
and the dump region addresses; the framing, handshake, encoding and checksum are common
code.

\section{Dump regions differ}
The dump region addresses are per-product. The SX-KN5000 and SX-KN1500 share eleven of
their thirteen region and length pairs with each other, but the SX-WSA1R shares none of
its eight with either --- its categories are its own.

\begin{caution}{The SX-WSA1R does not accept the GS-style messages its siblings do}
The SX-KN5000 and SX-KN1500 additionally accept a small set of messages in the
\bytes{42 12} parameter form, including the sequence \bytes{42 12 40 00 7F 00 41} used
elsewhere in the industry as a reset. \textbf{The SX-WSA1R accepts none of them.} Software
that resets a KN-series instrument that way and assumes the SX-WSA1R will follow will find
it does not respond.
\end{caution}
